\documentclass[11pt,letterpaper]{article}

% ------------------------------------------------------------------
% Paquetes
% ------------------------------------------------------------------
\usepackage[utf8]{inputenc}
\usepackage[T1]{fontenc}
\usepackage{lmodern}
\usepackage[spanish,es-tabla,es-noshorthands]{babel}
\usepackage{amsmath,amssymb}
\usepackage[margin=2.3cm,headheight=14pt]{geometry}
\usepackage{booktabs,array,tabularx,multirow,longtable}
\usepackage[table]{xcolor}
\usepackage{tikz}
\usetikzlibrary{positioning,arrows.meta,calc,fit,backgrounds,shapes.geometric}
\usepackage{listings}
\usepackage{siunitx}
\sisetup{output-decimal-marker={,}}
\usepackage{enumitem}
\usepackage{fancyhdr}
\usepackage[most]{tcolorbox}
\usepackage{caption}
\usepackage{float}
\usepackage[hidelinks]{hyperref}

% ------------------------------------------------------------------
% Datos del documento (editar aqui)
% ------------------------------------------------------------------
\newcommand{\universidad}{Universidad Tecnol\'ogica de Pereira}
\newcommand{\facultad}{Facultad de Ingenier\'ias}
\newcommand{\materia}{Arquitectura de Computadores}
\newcommand{\estudiante}{Juan Pablo Loaiza Bedoya}
% Compa\~neros de equipo (uno por l\'inea, separados con \\):
\newcommand*{\companeros}{Javier Parra Gonzales \\ Estiven Mu\~noz Montoya}
% Para poner el docente, escribir su nombre dentro de las llaves:
\newcommand*{\docente}{}
\newcommand{\fecha}{16 de noviembre de 2026}
\newcommand{\numtaller}{3}
\newcommand{\titulotaller}{El hardware se protege solo}

% ------------------------------------------------------------------
% Estilo
% ------------------------------------------------------------------
\definecolor{datos}{RGB}{31,94,168}
\definecolor{control}{RGB}{190,40,40}
\definecolor{reloj}{RGB}{30,130,60}
\definecolor{nuevo}{RGB}{200,110,0}
\definecolor{gris}{RGB}{245,245,245}
\definecolor{comentario}{RGB}{90,120,90}

\pagestyle{fancy}
\fancyhf{}
\fancyhead[L]{\small \materia{} -- Taller \numtaller}
\fancyhead[R]{\small UTP}
\fancyfoot[C]{\small \thepage}

\setlist{itemsep=2pt,topsep=4pt}
\setlength{\parskip}{4pt}
\setlength{\parindent}{0pt}
\renewcommand{\arraystretch}{1.15}

\lstdefinelanguage{SV}{
  morekeywords={module,endmodule,input,output,logic,parameter,localparam,int,
    assign,always_comb,always_ff,always,posedge,negedge,if,else,case,endcase,
    default,begin,end,import,function,endfunction,automatic,package,endpackage,
    integer,initial,for,string,task,endtask},
  sensitive=true,
  morecomment=[l]{//},
  morestring=[b]",
}
\lstdefinelanguage{RVASM}{
  morekeywords={addi,add,sub,and,or,xor,slt,sltu,lw,lb,lbu,lh,lhu,sw,sb,sh,
    beq,bne,blt,bge,bltu,bgeu,jal,jalr,lui,auipc,nop},
  sensitive=true,
  morecomment=[l]{\#},
}
\lstset{
  basicstyle=\ttfamily\footnotesize,
  keywordstyle=\color{datos}\bfseries,
  commentstyle=\color{comentario}\itshape,
  stringstyle=\color{control},
  backgroundcolor=\color{gris},
  frame=single, rulecolor=\color{black!25},
  columns=fullflexible, keepspaces=true,
  breaklines=true, breakatwhitespace=false,
  showstringspaces=false, tabsize=4,
  numbers=left, numberstyle=\tiny\color{black!50}, numbersep=6pt,
  xleftmargin=14pt, upquote=true,
  aboveskip=6pt, belowskip=6pt,
  literate={á}{{\'a}}1 {é}{{\'e}}1 {í}{{\'i}}1 {ó}{{\'o}}1 {ú}{{\'u}}1
           {Á}{{\'A}}1 {É}{{\'E}}1 {Í}{{\'I}}1 {Ó}{{\'O}}1 {Ú}{{\'U}}1
           {ñ}{{\~n}}1 {Ñ}{{\~N}}1 {ü}{{\"u}}1,
}

\tcbset{
  respuesta/.style={colback=green!4,colframe=reloj!70!black,boxrule=0.8pt,
    arc=2pt,left=6pt,right=6pt,top=4pt,bottom=4pt,fonttitle=\bfseries\small},
  enunciado/.style={colback=black!3,colframe=black!35,boxrule=0.5pt,arc=2pt,
    left=6pt,right=6pt,top=3pt,bottom=3pt,fonttitle=\bfseries\small},
  revisar/.style={colback=orange!6,colframe=nuevo,boxrule=0.8pt,arc=2pt,
    left=6pt,right=6pt,top=3pt,bottom=3pt,fonttitle=\bfseries\small},
}
\newtcolorbox{respuesta}[1][Resultado]{respuesta,title={#1}}
\newtcolorbox{enunciado}[1][Enunciado]{enunciado,title={#1}}
\newtcolorbox{supuesto}[1][Decisi\'on de dise\~no]{revisar,title={#1}}

\newcommand{\cod}[1]{\texttt{#1}}
\newcommand{\sig}[1]{\textcolor{control}{\texttt{#1}}}

% ------------------------------------------------------------------
\begin{document}

\begin{titlepage}
  \centering
  \vspace*{2.0cm}
  {\Huge\bfseries Taller \numtaller\par}
  \vspace{0.5cm}
  {\LARGE \titulotaller\par}
  \vfill
  {\large Presentado por:\par}
  \vspace{0.2cm}
  {\large\bfseries \estudiante\par}
  \ifx\companeros\empty\else
    \vspace{0.2cm}{\large\bfseries \companeros\par}
  \fi
  \ifx\docente\empty\else
    \vspace{0.8cm}
    {\large Docente:\par}\vspace{0.2cm}{\large \docente\par}
  \fi
  \vspace{2.0cm}
  {\large \materia\par}
  \vspace{0.3cm}
  {\Large\bfseries \universidad\par}
  \vspace{0.2cm}
  {\large \facultad\par}
  \vspace{0.2cm}
  {\large Ingenier\'ia de Sistemas y Computaci\'on\par}
  \vspace{0.8cm}
  {\large Pereira, \fecha\par}
\end{titlepage}

\tableofcontents
\newpage

% =====================================================================
\section{Introducci\'on}
% =====================================================================
\begin{enunciado}[Enunciado (resumen del PDF \emph{Taller 3 --- El hardware se protege solo})]
Partir del procesador monociclo del Taller 2 (el propio, que ya pasa su regresi\'on), cortarlo en cinco
etapas y a\~nadirle las dos unidades que hacen que un procesador segmentado funcione sin ayuda del
programador. Cuatro fases: (1) un registro de segmentaci\'on parametrizable con carga, habilitaci\'on y
borrado s\'incrono, y el corte del monociclo en cinco etapas; (2) la unidad de \emph{forwarding};
(3) la unidad de detecci\'on de riesgos; (4) el procesador segmentado cableado. Se entrega el
repositorio, los programas de prueba y un archivo corto en texto con: qu\'e situaciones se probaron y con
qu\'e programa, qu\'e sali\'o mal y c\'omo se encontr\'o, y cu\'antos ciclos tarda un programa en el
segmentado frente al monociclo. \textbf{Entrega: s\'abado 5 de diciembre, 23:59.}
\end{enunciado}

La gu\'ia describe el comportamiento de cada pieza sin dar c\'odigo; el trabajo consiste en traducir esa
descripci\'on a HDL. El informe sigue la misma estructura: para cada pieza se resume lo que pide la
gu\'ia, se explica la soluci\'on adoptada y se presenta el m\'odulo que la implementa. Al final se
presentan las pruebas realizadas y sus resultados.

\textbf{Metodolog\'ia de verificaci\'on.} El dise\~no se simul\'o con Icarus Verilog 12 (\cod{-g2012 -Wall}). Cada
programa corre a la vez en el monociclo del Taller~2, que hace de \emph{modelo de referencia}, y en el
segmentado, y se exige que el estado final (32 registros y 64 palabras de memoria) sea id\'entico.

\begin{table}[!htbp]
\centering\small
\begin{tabular}{@{}l l l@{}}
\toprule
\textbf{Archivo} & \textbf{Fase} & \textbf{Qu\'e es} \\
\midrule
\cod{rtl/pipe\_reg.sv}         & 1 & registro de segmentaci\'on: \cod{en}, \cod{clr}, \cod{rst} \\
\cod{rtl/forwarding\_unit.sv}  & 2 & unidad de adelantamiento (vive en EX) \\
\cod{rtl/hazard\_unit.sv}      & 3 & unidad de detecci\'on de riesgos (vive en ID) \\
\cod{rtl/riscv\_segmentado.sv} & 1 y 4 & el corte en cinco etapas y el cableado completo \\
\cod{tb/tb\_unidades.sv}       & pruebas & las tres piezas por separado (tablas de las secciones 6 y 12) \\
\cod{tb/tb\_pipeline.sv}       & pruebas & 11 programas: segmentado contra monociclo, ciclos \\
\cod{prog/*.s}                 & pruebas & 8 programas dirigidos nuevos + los 3 del Taller 2 \\
\cod{INFORME.txt}              & entrega & el archivo corto en texto que pide la gu\'ia \\
\bottomrule
\end{tabular}
\caption{Lo que se a\~nadi\'o al repositorio del Taller 2. Ning\'un m\'odulo del Taller 2 se modific\'o.}
\end{table}

% =====================================================================
\section{Fase 1: de monociclo a cinco etapas}
% =====================================================================

\subsection{D\'onde se corta}
\begin{center}\small
\begin{tabular}{@{}c l l@{}}
\toprule
\textbf{Etapa} & \textbf{Qu\'e hace} & \textbf{M\'odulos del Taller 2 que usa} \\
\midrule
IF  & busca la instrucci\'on & \cod{instruction\_memory}, \cod{pc}, sumador de $PC+4$, \cod{mux\_pc} \\
ID  & decodifica y lee el banco & \cod{control\_unit}, \cod{register\_file}, \cod{imm\_gen} \\
EX  & calcula & ALU, sumador del destino, \cod{mux\_jalr}, multiplexores de operando \\
MEM & lee o escribe memoria & \cod{data\_memory} \\
WB  & escribe en el banco & multiplexor de \cod{ResultSrc} \\
\bottomrule
\end{tabular}
\end{center}

\subsection{El registro de segmentaci\'on (\cod{pipe\_reg.sv})}
\begin{enunciado}
Un registro de ancho parametrizable, con carga s\'incrona, habilitaci\'on y borrado s\'incrono. Los tres
puertos hacen falta: sin habilitaci\'on no se puede parar; sin borrado no se puede meter una burbuja.
\end{enunciado}
Un solo m\'odulo sirve para los cuatro registros: cada uno empaqueta sus campos en un vector
\cod{d}/\cod{q} del ancho necesario. Prioridad: \cod{rst} y \cod{clr} ganan sobre \cod{en}. Un
registro puesto a cero \emph{es} una instrucci\'on muerta: \cod{RegWrite = MemWrite = Branch = Jump = 0}
y \cod{rd = x0}. Esa es la burbuja: un \cod{nop} que nadie escribi\'o y que no ocupa memoria.

\begin{lstlisting}[language=SV,title={\small\ttfamily rtl/pipe\_reg.sv}]
module pipe_reg #(
    parameter int W = 32
)(
    input  logic         clk,
    input  logic         rst,
    input  logic         en,
    input  logic         clr,
    input  logic [W-1:0] d,
    output logic [W-1:0] q
);

    always_ff @(posedge clk) begin
        if (rst || clr) q <= '0;
        else if (en)    q <= d;
    end

endmodule
\end{lstlisting}


\subsection{Qu\'e guarda cada registro, y cu\'antos bits}
La regla de la gu\'ia: \emph{lo que no se guarda, se pierde}. La tabla~\ref{tab:anchos} lista los campos
de cada registro tal como quedaron en \cod{riscv\_segmentado.sv} y verifica el ancho declarado en
cada instancia.

\begin{table}[!htbp]
\centering\small
\begin{tabularx}{\textwidth}{@{}l X r@{}}
\toprule
\textbf{Registro} & \textbf{Campos} & \textbf{Ancho} \\
\midrule
IF/ID  & \cod{Instr, PC, PC+4} & $3\cdot 32 = \mathbf{96}$ \\
ID/EX  & \cod{RD1, RD2, PC, PC+4, ImmExt} ($5\cdot 32 = 160$); \cod{rs1, rs2, rd} ($3\cdot 5 = 15$);
         \cod{funct3} (3); control: \cod{RegWrite}~1, \cod{ResultSrc}~2, \cod{MemWrite}~1, \cod{Jump}~1,
         \cod{Branch}~1, \cod{Jalr}~1, \cod{ALUControl}~4, \cod{ALUSrcA}~1, \cod{ALUSrcB}~1 (13)
       & $160+15+3+13 = \mathbf{191}$ \\
EX/MEM & resultado definitivo, dato a escribir, \cod{PC+4}, destino del salto ($4\cdot 32 = 128$);
         \cod{rd} (5); \cod{funct3} (3); \cod{RegWrite}~1, \cod{ResultSrc}~2, \cod{MemWrite}~1 (4)
       & $128+5+3+4 = \mathbf{140}$ \\
MEM/WB & dato le\'ido, resultado, \cod{PC+4} ($3\cdot 32 = 96$); \cod{rd} (5); \cod{RegWrite}~1,
         \cod{ResultSrc}~2 (3) & $96+5+3 = \mathbf{104}$ \\
\bottomrule
\end{tabularx}
\caption{Contenido y ancho de los cuatro registros de segmentaci\'on.}
\label{tab:anchos}
\end{table}

Tres campos que sorprenden, todos presentes:
\begin{itemize}
\item \textbf{\cod{PC} hasta ID/EX}, porque el destino $PC+\mathit{imm}$ se calcula en EX con el PC
  \emph{de esa instrucci\'on}, no con el de IF, que ya avanz\'o dos posiciones.
\item \textbf{\cod{PC+4} hasta MEM/WB}, porque \cod{jal} y \cod{jalr} escriben la direcci\'on de retorno.
\item \textbf{\cod{rs1} y \cod{rs2} hasta EX}, porque la unidad de forwarding compara
  \emph{n\'umeros} de registro, no valores.
\end{itemize}

\begin{supuesto}[Decisi\'on de dise\~no: campos adicionales en los registros de segmentaci\'on]
La tabla de la secci\'on~2 de la gu\'ia no incluye tres se\~nales que el procesador del curso necesita; por eso se a\~nadieron:
\begin{itemize}
\item \cod{funct3} viaja por ID/EX y EX/MEM. Hace falta en EX, para decidir el salto (\cod{beq},
  \cod{bne}, \ldots), y en MEM, porque decide el ancho del acceso a memoria (\cod{lb}/\cod{lh}/\cod{lw}).
  En el monociclo iba directo de la instrucci\'on a la memoria, pero en MEM la instrucci\'on ya no
  est\'a.
\item \cod{ALUSrcA} y \cod{Jalr} viajan como se\~nales de EX: la primera la usa \cod{auipc} y la
  segunda elige el destino de \cod{jalr}.
\end{itemize}
\end{supuesto}

\subsection{Las se\~nales de control y la decisi\'on del salto}
Todas las se\~nales de control se generan en ID con la misma \cod{control\_unit} del Taller~2 y viajan
con la instrucci\'on hasta la etapa en la que se consumen: \sig{ALUSrcA}, \sig{ALUSrc},
\sig{ALUControl}, \sig{Branch}, \sig{Jump} y \sig{Jalr} en EX; \sig{MemWrite} en MEM;
\sig{RegWrite} y \sig{ResultSrc} en WB.

\begin{supuesto}[Decisi\'on de dise\~no: d\'onde se decide \cod{PCSrc}]
En el monociclo, \cod{control\_unit} calcula \cod{pc\_src} a partir de las banderas de la ALU. En el
segmentado las banderas existen en EX, pero la unidad de control vive en ID. Para no modificar el
m\'odulo (la gu\'ia pide reutilizarlo ``sin un cambio''), se instancia en ID con las banderas a cero y
su salida \cod{pc\_src} sin conectar, y en EX se repite su bloque~3 (la tabla de \cod{funct3} contra
\cod{zero}/\cod{less}/\cod{less\_u}) sobre \cod{BranchE}, \cod{funct3E} y las banderas de la ALU:
$\boxed{\mathit{PCSrcE} = \mathit{BranchTakenE} \vee \mathit{JumpE}}$. Las banderas se calculan con
los operandos \emph{ya adelantados}, as\'i que un salto que depende de la instrucci\'on anterior compara
el valor bueno.
\end{supuesto}

\begin{supuesto}[Decisi\'on de dise\~no: escritura del banco en el primer medio ciclo]
La gu\'ia dice que la distancia 3 ya est\'a resuelta porque ``el banco de registros del curso escribe en
el primer medio ciclo y lee en el segundo''. El \cod{register\_file.sv} del Taller~2, en cambio,
escribe en el flanco de \emph{subida}. Sin tocar el m\'odulo, se le conecta el reloj invertido
(\cod{.clk(\textasciitilde clk)}): escribe en el flanco de bajada, a mitad del ciclo, y la lectura
combinacional de ID ve el valor nuevo antes del siguiente flanco de subida. Se comprob\'o que hace
falta: con el reloj sin invertir fallan \cod{completo.s}, \cod{doble.s} y \cod{lui\_jal.s}
(mutante M12, secci\'on~\ref{sec:mutantes}). Otra alternativa v\'alida habr\'ia sido un desv\'io interno
dentro del banco, pero exig\'ia modificar un m\'odulo entregado por el curso.
\end{supuesto}

% =====================================================================
\section{Fase 2: la unidad de forwarding}
% =====================================================================
\begin{enunciado}
Combinacional pura, vive en EX. Entradas: \cod{rs1E}, \cod{rs2E} (de ID/EX), \cod{rdM}, \cod{RegWriteM}
(de EX/MEM), \cod{rdW}, \cod{RegWriteW} (de MEM/WB). Salidas: \cod{ForwardAE} y \cod{ForwardBE}, de dos
bits: \cod{00} el banco, \cod{10} EX/MEM (distancia 1), \cod{01} MEM/WB (distancia 2). Regla: si el
registro que voy a leer coincide con el que va a escribir una instrucci\'on que llevo delante, tomo su
valor de donde est\'e; y si coinciden las dos, me quedo con la m\'as cercana.
\end{enunciado}

\subsection{La l\'ogica}
Se escribe una vez para cada operando: la gu\'ia advierte que implementar solo \cod{ForwardAE} es el
error m\'as frecuente, y es silencioso. Para el operando A:
\[
\mathit{ForwardAE} =
\begin{cases}
10 & \text{si } \mathit{rs1E}\neq 0 \;\wedge\; \mathit{RegWriteM} \;\wedge\; \mathit{rs1E} = \mathit{rdM}\\[2pt]
01 & \text{si no, y } \mathit{rs1E}\neq 0 \;\wedge\; \mathit{RegWriteW} \;\wedge\; \mathit{rs1E} = \mathit{rdW}\\[2pt]
00 & \text{en otro caso}
\end{cases}
\]
y $\mathit{ForwardBE}$ es la misma expresi\'on con $\mathit{rs2E}$. Los tres detalles de la
secci\'on~7 est\'an en cada rama:
\begin{enumerate}
\item \textbf{$\mathit{rs}\neq 0$.} El \cod{nop} can\'onico (\cod{addi x0, x0, 0}) tiene \cod{rd = x0}.
  Sin esta condici\'on, el resultado que se tira a \cod{x0} se adelanta y \cod{x0} deja de valer cero.
\item \textbf{\cod{RegWrite}.} En un \cod{sw} o un \cod{beq}, los bits 11:7 son del inmediato. Por
  ejemplo, \cod{sw x10, 8(x0)} tiene $\mathit{instr}[11{:}7] = 01000_2 = 8$: sin mirar \cod{RegWrite},
  se anunciar\'ia como una escritura a \cod{x8} y adelantar\'ia la direcci\'on (8).
\item \textbf{El orden es la prioridad.} EX/MEM se pregunta primero porque lleva el valor m\'as
  reciente (doble hazard, tabla~\ref{tab:traza_doble}).
\end{enumerate}

\subsection{Qu\'e valor se adelanta: el mux de \cod{ResultSrc} en EX}
Adelantar el resultado de la ALU no basta: \cod{jal}/\cod{jalr} escriben $PC+4$ y \cod{lui} escribe
el inmediato, valores que nunca pasan por la ALU. Por eso se a\~nadi\'o en EX una copia del mux de
\cod{ResultSrc} (\cod{mux\_res\_ex}) para los tres casos que ya tienen su valor en esa etapa:

\begin{center}\small
\begin{tabular}{@{}c l l@{}}
\toprule
\cod{ResultSrcE} & instrucciones & lo que entra a EX/MEM como resultado \\
\midrule
00 & \cod{add}, \cod{addi}, \cod{auipc}, \ldots & \cod{ALUResultE} \\
01 & cargas & \cod{ALUResultE} (la \emph{direcci\'on}: el dato a\'un no existe) \\
10 & \cod{jal}, \cod{jalr} & \cod{PCPlus4E} \\
11 & \cod{lui} & \cod{ImmExtE} \\
\bottomrule
\end{tabular}
\end{center}
As\'i, \cod{ALUResultM} es siempre el valor definitivo de \cod{rd}, salvo en las cargas, que es
justamente el caso que obliga a parar (fase~3). El mux de WB sigue eligiendo el dato le\'ido cuando
\cod{ResultSrcW = 01}.

Las entradas de los dos multiplexores de tres v\'ias (un \cod{mux4} con la rama \cod{11} sin usar) son
\cod{RD1E}/\cod{RD2E}, \cod{ResultW} y \cod{ALUResultM}. La salida del mux B es tambi\'en el dato que
escribe un \emph{store}: un \cod{sw} que guarda un registro reci\'en calculado necesita el valor
adelantado, no el que ley\'o el banco (mutante M13).

\begin{lstlisting}[language=SV,title={\small\ttfamily rtl/forwarding\_unit.sv}]
module forwarding_unit (
    input  logic [4:0] rs1E,
    input  logic [4:0] rs2E,
    input  logic [4:0] rdM,
    input  logic [4:0] rdW,
    input  logic       RegWriteM,
    input  logic       RegWriteW,
    output logic [1:0] ForwardAE,
    output logic [1:0] ForwardBE
);

    always_comb begin
        // ---- Operando A ----
        if      ((rs1E != 5'd0) && RegWriteM && (rs1E == rdM)) ForwardAE = 2'b10;
        else if ((rs1E != 5'd0) && RegWriteW && (rs1E == rdW)) ForwardAE = 2'b01;
        else                                                   ForwardAE = 2'b00;

        // ---- Operando B: la misma tabla con rs2E ----
        if      ((rs2E != 5'd0) && RegWriteM && (rs2E == rdM)) ForwardBE = 2'b10;
        else if ((rs2E != 5'd0) && RegWriteW && (rs2E == rdW)) ForwardBE = 2'b01;
        else                                                   ForwardBE = 2'b00;
    end

endmodule
\end{lstlisting}


\subsection{La tabla de la secci\'on 6, comprobada}
La tabla de la gu\'ia es la especificaci\'on. \cod{tb\_unidades} la aplica fila por fila, primero al
operando A y despu\'es al B, y comprueba en cada caso que el otro operando se queda en \cod{00}. Se
a\~nadieron dos filas extra.

\begin{table}[!htbp]
\centering\footnotesize\setlength{\tabcolsep}{4pt}
\begin{tabular}{@{}c c c c c c l c@{}}
\toprule
\textbf{rsE} & \textbf{rdM} & \textbf{RegWriteM} & \textbf{rdW} & \textbf{RegWriteW} & \textbf{Salida} &
\textbf{Situaci\'on} & \textbf{A y B} \\
\midrule
5  & 7  & 1 & 8 & 1 & 00 & nadie escribe ese registro & \checkmark \\
5  & 5  & 1 & 8 & 1 & 10 & la de delante escribe \cod{x5} & \checkmark \\
5  & 7  & 1 & 5 & 1 & 01 & la de dos por delante escribe \cod{x5} & \checkmark \\
5  & 5  & 1 & 5 & 1 & 10 & las dos: gana la m\'as cercana & \checkmark \\
5  & 5  & 0 & 8 & 1 & 00 & coincide, pero no escribe (un \cod{sw}) & \checkmark \\
5  & 7  & 1 & 5 & 0 & 00 & coincide, pero no escribe (un branch) & \checkmark \\
0  & 0  & 1 & 0 & 1 & 00 & el operando es \cod{x0}: nunca & \checkmark \\
31 & 31 & 1 & 8 & 1 & 10 & \cod{x31}: se comparan los cinco bits & \checkmark \\
\midrule
5  & 5  & 0 & 5 & 1 & 01 & \emph{extra:} la cercana no escribe, la lejana s\'i & \checkmark \\
0  & 7  & 1 & 0 & 1 & 00 & \emph{extra:} \cod{x0} solo en MEM/WB & \checkmark \\
\bottomrule
\end{tabular}
\caption{Comportamiento de la unidad de forwarding (40 comprobaciones en \cod{tb\_unidades}).}
\end{table}

\begin{table}[!htbp]
\centering\scriptsize
\setlength{\tabcolsep}{2.2pt}
\begin{tabular}{@{}lcccccccccccc@{}}
\toprule
\textbf{Instrucci\'on} & \textbf{c1} & \textbf{c2} & \textbf{c3} & \textbf{c4} & \textbf{c5} & \textbf{c6} & \textbf{c7} & \textbf{c8} & \textbf{c9} & \textbf{c10} & \textbf{c11} & \textbf{c12} \\
\midrule
\texttt{addi x1, x0, 1} & IF & ID & EX & MEM & WB &  &  &  &  &  &  &  \\
\texttt{addi x2, x0, 2} &  & IF & ID & EX & MEM & WB &  &  &  &  &  &  \\
\texttt{addi x3, x0, 30} &  &  & IF & ID & EX & MEM & WB &  &  &  &  &  \\
\texttt{addi x4, x0, 40} &  &  &  & IF & ID & EX & MEM & WB &  &  &  &  \\
\texttt{add x5, x1, x2} &  &  &  &  & IF & ID & EX & MEM & WB &  &  &  \\
\texttt{add x5, x3, x4} &  &  &  &  &  & IF & ID & EX & MEM & WB &  &  \\
\texttt{add x6, x5, x1} &  &  &  &  &  &  & IF & ID & EX & MEM & WB &  \\
\texttt{add x7, x1, x2} &  &  &  &  &  &  &  & IF & ID & EX & MEM & WB \\
\bottomrule
\end{tabular}
\caption{Traza real de \cod{doble.s} (doble hazard).}
\label{tab:traza_doble}
\end{table}

La tabla~\ref{tab:traza_doble} es la traza real de \cod{doble.s}. En el ciclo~c8, \cod{add x6, x5, x1}
est\'a en EX, el segundo \cod{add x5} en MEM y el primero en WB. Las dos condiciones se cumplen y gana
EX/MEM: $\mathit{x6} = 70 + 1 = 71$. Con la prioridad invertida saldr\'ia $3 + 1 = 4$, un valor
plausible y equivocado.

% =====================================================================
\section{Fase 3: la unidad de detecci\'on de riesgos}
% =====================================================================
\begin{enunciado}
Combinacional pura, vive en ID. Entradas: \cod{rs1D}, \cod{rs2D} (de IF/ID), \cod{rdE},
\cod{ResultSrcE} (de ID/EX) y \cod{PCSrcE} (de EX). Salidas: \cod{StallF} (habilitaci\'on del PC),
\cod{StallD} (habilitaci\'on de IF/ID), \cod{FlushD} (borrado de IF/ID) y \cod{FlushE} (borrado de
ID/EX). Dos trabajos: el \emph{load-use} y el salto tomado.
\end{enunciado}

\subsection{El load-use: la condici\'on completa}
El \'unico caso que el forwarding no resuelve es una carga seguida del uso de su \cod{rd}: el dato sale
de la memoria al final del mismo ciclo en el que la ALU lo necesita al principio. La soluci\'on es
retrasar un ciclo a la instrucci\'on dependiente; as\'i queda a distancia~2 y la cubre el camino
\cod{01} desde MEM/WB.
\[
\boxed{\;\mathit{lwStall} = (\mathit{ResultSrcE} = 01) \wedge \bigl((\mathit{rs1D} = \mathit{rdE})
\vee (\mathit{rs2D} = \mathit{rdE})\bigr)\;}
\]
\[
\boxed{\;\mathit{StallF} = \mathit{StallD} = \mathit{lwStall},\qquad
\mathit{FlushD} = \mathit{PCSrcE},\qquad
\mathit{FlushE} = \mathit{lwStall} \vee \mathit{PCSrcE}\;}
\]
\textbf{Por qu\'e se comparan los dos bits de \cod{ResultSrcE}.} Harris \& Harris comprueban solo el
bit~0 porque en su dise\~no \cod{ResultSrc} nunca vale \cod{11}. En este curso \cod{11} es el
inmediato de \cod{lui}, y \cod{11} tiene el bit~0 a uno: con la versi\'on del libro, un \cod{lui}
seguido del uso de su \cod{rd} provocar\'ia una parada innecesaria. Ning\'un programa lo delata, porque
una parada de m\'as no cambia el resultado. Solo lo detecta la prueba unitaria (mutante M11).

\textbf{Por qu\'e vive en ID y no en EX.} Cuando la dependiente llega a EX ya est\'a calculando con sus
operandos elegidos. Para no ejecutarla hay que detenerla antes, mientras est\'a en ID: la unidad de
forwarding corrige lo que ya va a pasar; la de riesgos impide que pase.

\textbf{Por qu\'e las tres se\~nales van juntas.} Si se para sin vaciar ID/EX, la copia que ya entr\'o a
EX se ejecuta y despu\'es se ejecuta la original: la instrucci\'on corre dos veces (mutante M7). Si se
vac\'ia ID/EX sin parar, IF/ID avanza y la instrucci\'on que estaba en ID se pierde (mutante M8).

\begin{lstlisting}[language=SV,title={\small\ttfamily rtl/hazard\_unit.sv}]
module hazard_unit (
    input  logic [4:0] rs1D,
    input  logic [4:0] rs2D,
    input  logic [4:0] rdE,
    input  logic [1:0] ResultSrcE,
    input  logic       PCSrcE,
    output logic       StallF,
    output logic       StallD,
    output logic       FlushD,
    output logic       FlushE
);

    import riscv_pkg::*;

    logic lwStall;

    assign lwStall = (ResultSrcE == RES_MEM) && ((rs1D == rdE) || (rs2D == rdE));

    assign StallF = lwStall;
    assign StallD = lwStall;
    assign FlushD = PCSrcE;
    assign FlushE = lwStall | PCSrcE;

endmodule
\end{lstlisting}


\begin{table}[!htbp]
\centering\footnotesize\setlength{\tabcolsep}{3.5pt}
\begin{tabular}{@{}c c c c c c c c l@{}}
\toprule
\textbf{rs1D} & \textbf{rs2D} & \textbf{rdE} & \textbf{ResultSrcE} & \textbf{PCSrcE} &
\textbf{Stall F/D} & \textbf{FlushD} & \textbf{FlushE} & \textbf{Situaci\'on} \\
\midrule
5 & 6 & 7 & 00 & 0 & 0 & 0 & 0 & nada que detectar \\
5 & 6 & 5 & 00 & 0 & 0 & 0 & 0 & depende de un \cod{add}: forwarding \\
5 & 6 & 5 & 01 & 0 & 1 & 0 & 1 & load-use por \cod{rs1} \\
5 & 6 & 6 & 01 & 0 & 1 & 0 & 1 & load-use por \cod{rs2} \\
5 & 6 & 7 & 01 & 0 & 0 & 0 & 0 & un \cod{lw}, pero a otro registro \\
5 & 6 & 5 & 10 & 0 & 0 & 0 & 0 & \cod{jal}: su $PC+4$ ya existe \\
5 & 6 & 5 & 11 & 0 & 0 & 0 & 0 & \cod{lui}: mirar solo el bit 0 parar\'ia de m\'as \\
5 & 6 & 7 & 00 & 1 & 0 & 1 & 1 & salto tomado \\
\bottomrule
\end{tabular}
\caption{Tabla de la secci\'on 12, verificada fila por fila (32 comprobaciones en \cod{tb\_unidades}).}
\end{table}

\begin{table}[!htbp]
\centering\scriptsize
\setlength{\tabcolsep}{2.2pt}
\begin{tabular}{@{}lcccccccccccccccc@{}}
\toprule
\textbf{Instrucci\'on} & \textbf{c1} & \textbf{c2} & \textbf{c3} & \textbf{c4} & \textbf{c5} & \textbf{c6} & \textbf{c7} & \textbf{c8} & \textbf{c9} & \textbf{c10} & \textbf{c11} & \textbf{c12} & \textbf{c13} & \textbf{c14} & \textbf{c15} & \textbf{c16} \\
\midrule
\texttt{addi x1, x0, 21} & IF & ID & EX & MEM & WB &  &  &  &  &  &  &  &  &  &  &  \\
\texttt{sw x1, 16(x0)} &  & IF & ID & EX & MEM & WB &  &  &  &  &  &  &  &  &  &  \\
\texttt{addi x2, x0, 4} &  &  & IF & ID & EX & MEM & WB &  &  &  &  &  &  &  &  &  \\
\texttt{sw x2, 20(x0)} &  &  &  & IF & ID & EX & MEM & WB &  &  &  &  &  &  &  &  \\
\texttt{lw x3, 16(x0)} &  &  &  &  & IF & ID & EX & MEM & WB &  &  &  &  &  &  &  \\
\texttt{add x4, x3, x0} &  &  &  &  &  & IF & ID & \textcolor{control}{\textbf{--}} & EX & MEM & WB &  &  &  &  &  \\
\texttt{lw x5, 20(x0)} &  &  &  &  &  &  & IF & \textcolor{control}{\textbf{--}} & ID & EX & MEM & WB &  &  &  &  \\
\texttt{add x6, x0, x5} &  &  &  &  &  &  &  &  & IF & ID & \textcolor{control}{\textbf{--}} & EX & MEM & WB &  &  \\
\texttt{lw x7, 16(x0)} &  &  &  &  &  &  &  &  &  & IF & \textcolor{control}{\textbf{--}} & ID & EX & MEM & WB &  \\
\texttt{addi x9, x0, 1} &  &  &  &  &  &  &  &  &  &  &  & IF & ID & EX & MEM & WB \\
\texttt{add x8, x7, x7} &  &  &  &  &  &  &  &  &  &  &  &  & IF & ID & EX & MEM \\
\texttt{lw x12, 16(x0)} &  &  &  &  &  &  &  &  &  &  &  &  &  & IF & ID & EX \\
\bottomrule
\end{tabular}
\caption{Traza real de \cod{lw\_use.s}: tres paradas de un ciclo.}
\label{tab:traza_lw}
\end{table}

La tabla~\ref{tab:traza_lw} es la traza real de \cod{lw\_use.s} (las marcas \textcolor{control}{\textbf{--}}
son los ciclos perdidos). Coincide con la de la secci\'on~10 de la gu\'ia: el \cod{add x4, x3, x0}
se queda en ID un ciclo, y la instrucci\'on de detr\'as (\cod{lw x5}) tambi\'en se queda quieta en IF,
porque se congela el frente entero. En cambio, \cod{lw x7}, seguido de \cod{addi x9} (independiente)
y de \cod{add x8, x7, x7}, no para: la distancia ya es~2. En total, el programa tiene exactamente 3
paradas, y el testbench lo exige.

\subsection{El salto tomado}
La condici\'on se resuelve en EX. Para entonces, las dos instrucciones de detr\'as ya entraron;
\cod{FlushD} vac\'ia IF/ID y \cod{FlushE} vac\'ia ID/EX. Las dos causas de \cod{FlushE} no pueden
coincidir: una exige una carga en EX y la otra un salto en EX.

\begin{table}[!htbp]
\centering\scriptsize
\setlength{\tabcolsep}{2.2pt}
\begin{tabular}{@{}lcccccccccccc@{}}
\toprule
\textbf{Instrucci\'on} & \textbf{c1} & \textbf{c2} & \textbf{c3} & \textbf{c4} & \textbf{c5} & \textbf{c6} & \textbf{c7} & \textbf{c8} & \textbf{c9} & \textbf{c10} & \textbf{c11} & \textbf{c12} \\
\midrule
\texttt{addi x1, x0, 3} & IF & ID & EX & MEM & WB &  &  &  &  &  &  &  \\
\texttt{addi x2, x0, 3} &  & IF & ID & EX & MEM & WB &  &  &  &  &  &  \\
\texttt{add x3, x1, x0} &  &  & IF & ID & EX & MEM & WB &  &  &  &  &  \\
\texttt{beq x3, x2, destino} &  &  &  & IF & ID & EX & MEM & WB &  &  &  &  \\
\textcolor{control}{\texttt{addi x30, x0, 99}\ (borrada)} &  &  &  &  & IF & ID &  &  &  &  &  &  \\
\textcolor{control}{\texttt{sw x1, 0(x0)}\ (borrada)} &  &  &  &  &  & IF &  &  &  &  &  &  \\
\texttt{addi x4, x0, 7} &  &  &  &  &  &  & IF & ID & EX & MEM & WB &  \\
\texttt{jal x0, sig} &  &  &  &  &  &  &  & IF & ID & EX & MEM & WB \\
\textcolor{control}{\texttt{addi x31, x0, 99}\ (borrada)} &  &  &  &  &  &  &  &  & IF & ID &  &  \\
\textcolor{control}{\texttt{sw x1, 4(x0)}\ (borrada)} &  &  &  &  &  &  &  &  &  & IF &  &  \\
\bottomrule
\end{tabular}
\caption{Traza real de \cod{salto.s}: dos saltos tomados, cuatro instrucciones borradas.}
\label{tab:traza_salto}
\end{table}

En la tabla~\ref{tab:traza_salto}, el \cod{beq} est\'a en EX en el ciclo~c6 y salta: \cod{addi x30} y
\cod{sw x1, 0(x0)} desaparecen sin escribir nada, y en c7 entra la instrucci\'on del destino. Lo mismo
pasa con \cod{jal}. La penalizaci\'on es de 2 ciclos por salto tomado.

% =====================================================================
\section{Fase 4: el procesador segmentado}
% =====================================================================

\begin{figure}[!htbp]
\centering
\resizebox{\textwidth}{!}{%
\begin{tikzpicture}[
  >=Stealth, font=\footnotesize,
  caja/.style={draw, thick, rounded corners=2pt, align=center, fill=white},
  dada/.style={caja, fill=black!7},
  nueva/.style={caja, draw=nuevo, fill=orange!8, very thick},
  preg/.style={draw, thick, fill=black!18, minimum width=0.38cm, minimum height=7.2cm},
  mux/.style={draw, thick, trapezium, trapezium angle=70, shape border rotate=270,
              minimum width=1.25cm, minimum height=0.5cm, inner sep=1pt, fill=white},
  d/.style={->, thick, datos},
  n/.style={->, thick, nuevo},
  c/.style={->, control, thick},
  et/.style={font=\tiny, datos},
  ec/.style={font=\tiny, control},
]
% registros de segmentacion
\node[preg] (r1) at (2.8,0.6) {};  \node[font=\tiny, above] at (r1.north) {IF/ID};
\node[preg] (r2) at (7.6,0.6) {};  \node[font=\tiny, above] at (r2.north) {ID/EX};
\node[preg] (r3) at (13.0,0.6) {}; \node[font=\tiny, above] at (r3.north) {EX/MEM};
\node[preg] (r4) at (16.4,0.6) {}; \node[font=\tiny, above] at (r4.north) {MEM/WB};
% etapas
\foreach \x/\t in {0.9/IF, 5.2/ID, 10.3/EX, 14.7/MEM, 17.6/WB}
  \node[font=\small\bfseries] at (\x,4.9) {\t};
% IF
\node[dada, minimum width=0.9cm, minimum height=1.2cm] (pc) at (0,0.6) {PC};
\node[dada, minimum width=1.2cm, minimum height=1.6cm] (im) at (1.6,0.6) {Mem.\\instr.};
\draw[d] (pc.east) -- (im.west);
\draw[d] (im.east) -- (r1.west |- im.east);
% ID
\node[caja, draw=control, minimum width=1.9cm, minimum height=0.8cm] (cu) at (5.2,3.3) {Unidad de\\control};
\node[dada, minimum width=1.9cm, minimum height=2.0cm] (rf) at (5.2,0.6) {Banco de\\registros};
\node[caja, minimum width=1.9cm, minimum height=0.7cm] (ig) at (5.2,-1.4) {Gen. inmed.};
\draw[d] (r1.east |- 0,0.6) -- (rf.west);
\draw[d] (3.4,0.6) |- (cu.west);
\draw[d] (3.4,0.6) |- (ig.west);
\draw[d] (rf.east |- 0,1.1) -- node[et, above]{\cod{RD1}} (r2.west |- 0,1.1);
\draw[d] (rf.east |- 0,0.1) -- node[et, above]{\cod{RD2}} (r2.west |- 0,0.1);
\draw[d] (ig.east) -- (r2.west |- ig.east);
\draw[->, control, dashed] (cu.east) -- (r2.west |- cu.east) node[pos=0.5, above, ec]{control};
% EX
\node[mux, draw=nuevo] (fa) at (8.6,1.4) {};
\node[mux, draw=nuevo] (fb) at (8.6,-0.2) {};
\node[mux] (sa) at (9.9,1.4) {};
\node[mux] (sb) at (9.9,-0.2) {};
\node[dada, minimum width=1.0cm, minimum height=2.4cm] (alu) at (11.3,0.6) {ALU};
\node[caja, minimum width=1.0cm, minimum height=0.6cm, draw=nuevo] (mre) at (12.15,-1.8) {\tiny mux\\[-2pt]\tiny ResSrc};
\draw[d] (r2.east |- 0,1.65) -- ($(fa.west)+(0,0.25)$);
\draw[d] (r2.east |- 0,0.05) -- ($(fb.west)+(0,0.25)$);
\draw[d] (fa.east) -- (sa.west |- fa.east);
\draw[d] (fb.east) -- (sb.west |- fb.east);
\draw[d] (sa.east) -- (alu.west |- sa.east);
\draw[d] (sb.east) -- (alu.west |- sb.east);
\draw[d] (alu.east) -- (12.15,0.6) -- (mre.north);
\draw[d] (mre.east) -- ++(0.3,0) |- (r3.west |- 0,-1.0) ;
\node[ec] at (8.6,2.2) {ForwardAE};
\node[ec, right] at (8.95,-0.95) {ForwardBE};
% MEM
\node[caja, minimum width=1.6cm, minimum height=2.0cm] (dm) at (14.7,0.6) {Memoria\\de datos};
\draw[d] (r3.east |- 0,0.6) -- (dm.west);
\draw[d] (dm.east) -- (r4.west |- dm.east);
% WB
\node[mux] (mr) at (17.4,0.6) {};
\draw[d] (r4.east |- 0,0.6) -- (mr.west);
% caminos de vuelta (forwarding)
\draw[n] (13.35,-1.0) -- (13.6,-1.0) -- (13.6,-2.55) -- (8.25,-2.55) -- (8.25,-0.45) -- ($(fb.west)+(0,-0.25)$);
\draw[n] (8.25,-2.55) -- (8.05,-2.55) -- (8.05,1.15) -- ($(fa.west)+(0,-0.25)$);
\node[font=\tiny, nuevo, below] at (11.0,-2.55) {\cod{ALUResultM} $\cdot$ de EX/MEM};
\draw[n] (mr.east) -- (18.1,0.6) -- (18.1,-3.15) -- (7.9,-3.15) -- (7.9,1.4);
\draw[n] (7.9,1.4) -- ($(fa.west)+(0,0)$);
\draw[n] (7.9,-0.2) -- ($(fb.west)+(0,0)$);
\fill[nuevo] (7.9,-0.2) circle (1.3pt);
\node[font=\tiny, nuevo, below] at (15.4,-3.15) {\cod{ResultW} $\cdot$ de MEM/WB (tambi\'en vuelve al banco)};
% unidades nuevas
\node[nueva, minimum width=3.6cm, minimum height=0.9cm] (hu) at (3.4,-4.4) {Unidad de detecci\'on de riesgos\\\tiny vive en ID};
\node[nueva, minimum width=3.6cm, minimum height=0.9cm] (fu) at (10.3,-4.4) {Unidad de forwarding\\\tiny vive en EX};
\node[ec, below] at (3.4,-4.95) {entradas: rs1D rs2D rdE ResultSrcE PCSrcE};
\node[ec, below] at (10.3,-4.95) {entradas: rs1E rs2E rdM rdW RegWriteM RegWriteW};
\draw[c] (hu.west) -| (pc.south) node[pos=0.75, left, ec]{StallF};
\draw[c] ($(hu.north)+(-0.6,0)$) -- (2.8,-3.0) node[pos=0.5, left, ec]{StallD, FlushD} -- (r1.south);
\draw[c] ($(hu.north)+(0.8,0)$) |- (7.6,-3.7) node[pos=0.75, below, ec]{FlushE} -- (r2.south);
\draw[c] ($(fu.north)+(-1.4,0)$) -- (8.9,-3.6) -- (8.9,-3.6);
\draw[c] (8.9,-3.6) -- (8.9,-3.4) -| (fb.south);
\node[ec, right] at (8.9,-3.45) {Forward A/B};
% salto
\draw[d] (alu.north) -- ++(0,0.9) node[et, above]{\cod{PCTargetE}, \cod{PCSrcE} $\to$ IF};
\end{tikzpicture}}
\caption{El procesador segmentado: cuatro registros de segmentaci\'on, la unidad de riesgos en ID y la
de forwarding en EX. En naranja, lo que a\~nade este taller: las dos unidades, los dos multiplexores de
tres entradas, la copia del mux de \cod{ResultSrc} en EX y los dos caminos de vuelta.}
\label{fig:pipeline}
\end{figure}

Aparte de las dos cajas nuevas, no cambi\'o nada respecto al monociclo: los mismos sumadores, el mismo
banco, las mismas dos memorias, la misma unidad de control y la misma ALU. Se a\~nadieron los cuatro
registros, dos bloques combinacionales, dos multiplexores de tres entradas, la copia del mux de
\cod{ResultSrc} en EX y los cables. Las se\~nales de parada y borrado llegan a estos puertos:

\begin{center}\small
\begin{tabular}{@{}l l l@{}}
\toprule
\textbf{Se\~nal} & \textbf{Puerto} & \textbf{Efecto} \\
\midrule
\sig{StallF} & \cod{registro\_pc.en = en \& \textasciitilde StallF} & el PC no avanza: se vuelve a buscar la misma instrucci\'on \\
\sig{StallD} & \cod{reg\_if\_id.en = en \& \textasciitilde StallD} & la dependiente se queda en ID \\
\sig{FlushD} & \cod{reg\_if\_id.clr} & mata la instrucci\'on que ven\'ia en IF \\
\sig{FlushE} & \cod{reg\_id\_ex.clr} & mete una burbuja en EX \\
\sig{ForwardAE} & \cod{mux\_fwd\_a.sel} & elige \cod{RD1E} / \cod{ResultW} / \cod{ALUResultM} \\
\sig{ForwardBE} & \cod{mux\_fwd\_b.sel} & elige \cod{RD2E} / \cod{ResultW} / \cod{ALUResultM} \\
\bottomrule
\end{tabular}
\end{center}

\begin{lstlisting}[language=SV,title={\small\ttfamily rtl/riscv\_segmentado.sv}]
module riscv_segmentado #(
    parameter int IMEM_DEPTH = 256,
    parameter int DMEM_DEPTH = 256,
    parameter     INIT_FILE  = ""
)(
    input logic clk,
    input logic rst,
    input logic en
);

    import riscv_pkg::*;

    // =========================================================================
    // Senales de las dos unidades nuevas
    // =========================================================================
    logic       StallF, StallD, FlushD, FlushE;
    logic [1:0] ForwardAE, ForwardBE;

    // =========================================================================
    // IF - busqueda
    // =========================================================================
    logic [31:0] PCF, PCPlus4F, PCNextF, InstrF;
    logic [31:0] PCTargetE;          // vuelve desde EX
    logic        PCSrcE;             // vuelve desde EX

    pc registro_pc (
        .clk     (clk),
        .rst     (rst),
        .en      (en & ~StallF),
        .pc_next (PCNextF),
        .pc      (PCF)
    );

    adder add_pc4 (
        .a (PCF),
        .b (32'd4),
        .y (PCPlus4F)
    );

    instruction_memory #(
        .DEPTH     (IMEM_DEPTH),
        .INIT_FILE (INIT_FILE)
    ) imem (
        .addr  (PCF),
        .instr (InstrF)
    );

    mux2 #(.W(32)) mux_pc (
        .d0  (PCPlus4F),
        .d1  (PCTargetE),
        .sel (PCSrcE),
        .y   (PCNextF)
    );

    // =========================================================================
    // IF/ID: la instruccion, el PC y PC+4                         (96 bits)
    // =========================================================================
    logic [31:0] InstrD, PCD, PCPlus4D;

    pipe_reg #(.W(96)) reg_if_id (
        .clk (clk),
        .rst (rst),
        .en  (en & ~StallD),
        .clr (FlushD),
        .d   ({InstrF, PCF, PCPlus4F}),
        .q   ({InstrD, PCD, PCPlus4D})
    );

    // =========================================================================
    // ID - decodificacion y lectura del banco
    // =========================================================================
    logic [6:0]  opcodeD;
    logic [2:0]  funct3D;
    logic        funct7b5D;
    logic [4:0]  rs1D, rs2D, rdD;

    assign opcodeD   = InstrD[6:0];
    assign rdD       = InstrD[11:7];
    assign funct3D   = InstrD[14:12];
    assign rs1D      = InstrD[19:15];
    assign rs2D      = InstrD[24:20];
    assign funct7b5D = InstrD[30];

    logic        RegWriteD, ALUSrcAD, ALUSrcBD, MemWriteD;
    logic        BranchD, JumpD, JalrD;
    logic [2:0]  ImmSrcD;
    logic [1:0]  ResultSrcD;
    logic [3:0]  ALUControlD;
    logic [31:0] RD1D, RD2D, ImmExtD;

    // Lo que se escribe en el banco, desde WB
    logic [31:0] ResultW;
    logic [4:0]  rdW;
    logic        RegWriteW;

    // La unidad de control es la del Taller 2, sin cambios. En ID solo se usan
    // sus senales "estaticas"; las banderas todavia no existen (se calculan en
    // EX), asi que entran a cero y su pc_src no se usa: el salto se decide en EX.
    control_unit control (
        .opcode      (opcodeD),
        .funct3      (funct3D),
        .funct7b5    (funct7b5D),
        .zero        (1'b0),
        .less        (1'b0),
        .less_u      (1'b0),
        .reg_write   (RegWriteD),
        .imm_src     (ImmSrcD),
        .alu_src_a   (ALUSrcAD),
        .alu_src_b   (ALUSrcBD),
        .mem_write   (MemWriteD),
        .result_src  (ResultSrcD),
        .branch      (BranchD),
        .jump        (JumpD),
        .jalr        (JalrD),
        .alu_control (ALUControlD),
        .pc_src      ()
    );

    // El banco escribe en el flanco de BAJADA (primer medio ciclo) y se lee
    // combinacionalmente despues: una escritura en WB y una lectura en ID del
    // mismo ciclo ya entregan el valor nuevo. Asi se resuelve la distancia 3
    // sin tocar register_file.sv: solo se le conecta el reloj invertido.
    register_file rf (
        .clk (~clk),
        .we  (RegWriteW),
        .rs1 (rs1D),
        .rs2 (rs2D),
        .rd  (rdW),
        .wd3 (ResultW),
        .rd1 (RD1D),
        .rd2 (RD2D)
    );

    imm_gen ig (
        .instr   (InstrD),
        .imm_src (ImmSrcD),
        .imm_ext (ImmExtD)
    );

    // Unidad de deteccion de riesgos: mira IF/ID (rs1D, rs2D) e ID/EX (rdE,
    // ResultSrcE). Vive en ID porque hay que detener la instruccion ANTES de
    // que llegue a EX.
    logic [4:0] rdE;
    logic [1:0] ResultSrcE;

    hazard_unit hu (
        .rs1D       (rs1D),
        .rs2D       (rs2D),
        .rdE        (rdE),
        .ResultSrcE (ResultSrcE),
        .PCSrcE     (PCSrcE),
        .StallF     (StallF),
        .StallD     (StallD),
        .FlushD     (FlushD),
        .FlushE     (FlushE)
    );

    // =========================================================================
    // ID/EX                                                       (191 bits)
    //   datos:   RD1, RD2, PC, PC+4, ImmExt                        5 x 32
    //   numeros: rs1, rs2, rd                                      3 x 5
    //   funct3                                                     3
    //   control: RegWrite, ResultSrc(2), MemWrite, Jump, Branch,
    //            Jalr, ALUControl(4), ALUSrcA, ALUSrcB             13
    // rs1 y rs2 viajan hasta EX: la unidad de forwarding compara NUMEROS.
    // funct3 viaja porque el salto se decide en EX y la memoria lo usa en MEM.
    // =========================================================================
    logic [31:0] RD1E, RD2E, PCE, PCPlus4E, ImmExtE;
    logic [4:0]  rs1E, rs2E;
    logic [2:0]  funct3E;
    logic        RegWriteE, MemWriteE, JumpE, BranchE, JalrE;
    logic        ALUSrcAE, ALUSrcBE;
    logic [3:0]  ALUControlE;

    pipe_reg #(.W(191)) reg_id_ex (
        .clk (clk),
        .rst (rst),
        .en  (en),
        .clr (FlushE),
        .d   ({RD1D, RD2D, PCD, PCPlus4D, ImmExtD,
               rs1D, rs2D, rdD, funct3D,
               RegWriteD, ResultSrcD, MemWriteD, JumpD, BranchD, JalrD,
               ALUControlD, ALUSrcAD, ALUSrcBD}),
        .q   ({RD1E, RD2E, PCE, PCPlus4E, ImmExtE,
               rs1E, rs2E, rdE, funct3E,
               RegWriteE, ResultSrcE, MemWriteE, JumpE, BranchE, JalrE,
               ALUControlE, ALUSrcAE, ALUSrcBE})
    );

    // =========================================================================
    // EX - ejecucion
    // =========================================================================
    logic [31:0] ALUResultM;         // vuelve desde EX/MEM (distancia 1)
    logic [4:0]  rdM;
    logic        RegWriteM;

    forwarding_unit fu (
        .rs1E      (rs1E),
        .rs2E      (rs2E),
        .rdM       (rdM),
        .rdW       (rdW),
        .RegWriteM (RegWriteM),
        .RegWriteW (RegWriteW),
        .ForwardAE (ForwardAE),
        .ForwardBE (ForwardBE)
    );

    // Los dos multiplexores de tres entradas (mux4 con la rama 11 sin uso)
    logic [31:0] FwdAE, WriteDataE;

    mux4 #(.W(32)) mux_fwd_a (
        .d00 (RD1E),
        .d01 (ResultW),
        .d10 (ALUResultM),
        .d11 (32'd0),
        .sel (ForwardAE),
        .y   (FwdAE)
    );

    // La salida del mux B es tambien el dato que un store escribe en memoria:
    // un sw que guarda un registro recien calculado necesita el valor bueno.
    mux4 #(.W(32)) mux_fwd_b (
        .d00 (RD2E),
        .d01 (ResultW),
        .d10 (ALUResultM),
        .d11 (32'd0),
        .sel (ForwardBE),
        .y   (WriteDataE)
    );

    // Despues del forwarding, los mismos dos mux2 del monociclo
    logic [31:0] SrcAE, SrcBE, ALUResultE;
    logic        zeroE, lessE, less_uE;

    mux2 #(.W(32)) mux_a (
        .d0  (FwdAE),
        .d1  (PCE),
        .sel (ALUSrcAE),
        .y   (SrcAE)
    );

    mux2 #(.W(32)) mux_b (
        .d0  (WriteDataE),
        .d1  (ImmExtE),
        .sel (ALUSrcBE),
        .y   (SrcBE)
    );

    alu alu_i (
        .src_a       (SrcAE),
        .src_b       (SrcBE),
        .alu_control (ALUControlE),
        .alu_result  (ALUResultE),
        .zero        (zeroE),
        .less        (lessE),
        .less_u      (less_uE)
    );

    // Destino del salto: PC de ESTA instruccion + inmediato, o rs1+imm (jalr)
    logic [31:0] PCPlusImmE;

    adder add_tgt (
        .a (PCE),
        .b (ImmExtE),
        .y (PCPlusImmE)
    );

    mux2 #(.W(32)) mux_jalr (
        .d0  (PCPlusImmE),
        .d1  ({ALUResultE[31:1], 1'b0}),
        .sel (JalrE),
        .y   (PCTargetE)
    );

    // Decision del salto en EX. Es exactamente el bloque 3 de control_unit,
    // aplicado a la instruccion que esta en EX con las banderas de la ALU,
    // que se calculan de los operandos YA adelantados.
    logic BranchTakenE;

    always_comb begin
        BranchTakenE = 1'b0;
        if (BranchE) begin
            case (funct3E)
                F3_BEQ : BranchTakenE =  zeroE;
                F3_BNE : BranchTakenE = ~zeroE;
                F3_BLT : BranchTakenE =  lessE;
                F3_BGE : BranchTakenE = ~lessE;
                F3_BLTU: BranchTakenE =  less_uE;
                F3_BGEU: BranchTakenE = ~less_uE;
                default: BranchTakenE = 1'b0;
            endcase
        end
    end

    assign PCSrcE = BranchTakenE | JumpE;

    // El mux de ResultSrc, adelantado a EX para los tres casos que ya tienen
    // su valor aqui (ALU, PC+4, ImmExt). Para un load (01) deja pasar la
    // direccion, que es lo que necesita la memoria en MEM. Asi EX/MEM lleva
    // siempre el valor definitivo de rd y el forwarding de un lui, un jal o
    // un auipc no adelanta basura.
    logic [31:0] ExResultE;

    mux4 #(.W(32)) mux_res_ex (
        .d00 (ALUResultE),
        .d01 (ALUResultE),
        .d10 (PCPlus4E),
        .d11 (ImmExtE),
        .sel (ResultSrcE),
        .y   (ExResultE)
    );

    // =========================================================================
    // EX/MEM                                                      (140 bits)
    //   resultado (definitivo), dato a escribir, PC+4, destino del salto,
    //   rd, funct3 y las senales de MEM y WB
    // =========================================================================
    logic [31:0] WriteDataM, PCPlus4M, PCTargetM;
    logic [2:0]  funct3M;
    logic [1:0]  ResultSrcM;
    logic        MemWriteM;

    pipe_reg #(.W(140)) reg_ex_mem (
        .clk (clk),
        .rst (rst),
        .en  (en),
        .clr (1'b0),
        .d   ({ExResultE,  WriteDataE, PCPlus4E, PCTargetE,
               rdE, funct3E, RegWriteE, ResultSrcE, MemWriteE}),
        .q   ({ALUResultM, WriteDataM, PCPlus4M, PCTargetM,
               rdM, funct3M, RegWriteM, ResultSrcM, MemWriteM})
    );

    // =========================================================================
    // MEM - memoria de datos
    // =========================================================================
    logic [31:0] ReadDataM;

    data_memory #(
        .DEPTH (DMEM_DEPTH)
    ) dmem (
        .clk    (clk),
        .we     (MemWriteM),
        .funct3 (funct3M),
        .addr   (ALUResultM),
        .wd     (WriteDataM),
        .rd     (ReadDataM)
    );

    // =========================================================================
    // MEM/WB                                                      (104 bits)
    //   dato leido, resultado, PC+4, rd y las senales de WB
    // =========================================================================
    logic [31:0] ReadDataW, ALUResultW, PCPlus4W;
    logic [1:0]  ResultSrcW;

    pipe_reg #(.W(104)) reg_mem_wb (
        .clk (clk),
        .rst (rst),
        .en  (en),
        .clr (1'b0),
        .d   ({ReadDataM, ALUResultM, PCPlus4M, rdM, RegWriteM, ResultSrcM}),
        .q   ({ReadDataW, ALUResultW, PCPlus4W, rdW, RegWriteW, ResultSrcW})
    );

    // =========================================================================
    // WB - escritura en el banco
    // =========================================================================
    // Las ramas 10 y 11 ya vienen resueltas desde EX dentro de ALUResultW; se
    // deja PC+4 en la 10 porque es el mismo valor y asi la tabla de ResultSrc
    // se lee igual que en el monociclo.
    mux4 #(.W(32)) mux_res (
        .d00 (ALUResultW),
        .d01 (ReadDataW),
        .d10 (PCPlus4W),
        .d11 (ALUResultW),
        .sel (ResultSrcW),
        .y   (ResultW)
    );

endmodule
\end{lstlisting}


% =====================================================================
\section{C\'omo se supo que est\'a bien}
% =====================================================================

\subsection{Los programas dirigidos (secci\'on 19)}
Un procesador segmentado mal construido pasa muchos programas; la prueba de humo del Taller~2 sigue en
verde con la unidad de forwarding a medias. Por eso se escribi\'o un programa para cada situaci\'on
de la secci\'on~19:

\begin{table}[!htbp]
\centering\small
\begin{tabularx}{\textwidth}{@{}X l X@{}}
\toprule
\textbf{Hay que provocar} & \textbf{Programa} & \textbf{Qu\'e se comprueba} \\
\midrule
Distancia 1 y 2, por A y por B & \cod{fwd\_dist.s} & con \cod{sub}, para que el orden importe:
  \cod{x3 = 4}, \cod{x4 = -1}, \cod{x6 = 1}, \cod{x7 = -1} \\
Dos escrituras seguidas y una lectura & \cod{doble.s} & \cod{x6 = x8 = 71} (al rev\'es saldr\'ia 4) \\
\cod{rd = x0} y lectura de \cod{x0} & \cod{x0.s} & \cod{x0 = 0}, \cod{x2 = 5}, \cod{x3 = 1}, \cod{x4 = 5} \\
\cod{sw} con desplazamiento = registro le\'ido & \cod{sw\_rd.s} & \cod{sw x10, 8(x0)}, luego leer \cod{x8}:
  \cod{x11 = 100} (no 8); tambi\'en con \cod{sw 12} y un \cod{bne} \\
\cod{lui} y uso a distancia 1 & \cod{lui\_jal.s} & tambi\'en \cod{jal}, \cod{auipc} y \cod{jalr}:
  \cod{x2 = 0x12345678}, \cod{x6 = 0x14} \\
\cod{lw} y uso, por \cod{rs1} y por \cod{rs2} & \cod{lw\_use.s} & \cod{x4 = 21}, \cod{x6 = 4}, y
  \emph{exactamente} 3 paradas \\
\cod{lw}, independiente, uso & \cod{lw\_use.s} & \cod{x8 = 42} sin parada (la cuenta de 3 lo exige) \\
Salto tomado con ranuras visibles & \cod{salto.s} & los seis condicionales tomados y no tomados, y
  \cod{jal}: \cod{x28..x31 = 0}, \cod{Mem[0..1] = 0} \\
\cod{lw} y uso con acumulador & \cod{acumulador.s} & \cod{x2 = 15}, \cod{x5 = 3}, \cod{Mem[5] = 15} \\
\bottomrule
\end{tabularx}
\caption{Situaciones de la secci\'on 19 y el programa que las provoca.}
\end{table}

Adem\'as de los valores puntuales, cada programa debe dejar el \emph{mismo} estado que el monociclo
(32 registros y 64 palabras). Resultado: \cod{tb\_unidades} \textbf{78/78} y \cod{tb\_pipeline}
\textbf{75/75}. La regresi\'on del Taller 2 tambi\'en sigue pasando (\textbf{86/86}).

\subsection{Defectos inyectados a prop\'osito}\label{sec:mutantes}
Que todo pase no demuestra que las pruebas sirvan. Para comprobarlo, se introdujo cada defecto de la
secci\'on~19 (uno a la vez, en una copia del repositorio) y se miró qu\'e prueba lo delata. Un defecto
que no hace fallar nada revela un hueco en las pruebas.

\begin{table}[!htbp]
\centering\footnotesize
\begin{tabularx}{\textwidth}{@{}l X c c X@{}}
\toprule
& \textbf{Defecto inyectado} & \textbf{Fallos unid.} & \textbf{Fallos sist.} & \textbf{Programas que lo delatan} \\
\midrule
M1 & Solo ForwardAE (ForwardBE siempre 00) & 5 & 36 & acumulador, bytes, completo, doble, fwd\_dist, humo, lui\_jal, lw\_use, salto, x0 \\
M2 & Prioridad invertida (MEM/WB antes que EX/MEM) & 1 & 4 & completo, doble \\
M3 & Sin la condicion rs $\neq$ x0 & 2 & 3 & x0 \\
M4 & Sin la condicion RegWrite & 2 & 2 & sw\_rd \\
M5 & EX/MEM solo lleva el resultado de la ALU & 0 & 8 & bytes, completo, lui\_jal \\
M6 & Sin deteccion de load-use & 6 & 9 & acumulador, lw\_use \\
M7 & Parar sin vaciar ID/EX & 2 & 3 & acumulador \\
M8 & Vaciar ID/EX sin parar & 4 & 9 & acumulador, lw\_use \\
M9 & Salto tomado sin FlushD & 1 & 7 & completo, salto \\
M10 & Salto tomado sin FlushE & 1 & 8 & completo, lui\_jal, salto \\
M11 & Comparar solo \cod{ResultSrcE[0]} (version del libro) & 3 & 0 & - \\
M12 & Banco con el reloj sin invertir (escribe en subida) & 0 & 8 & completo, doble, lui\_jal \\
M13 & Store con RD2E en vez del dato adelantado & 0 & 15 & acumulador, bytes, humo, lw\_use \\
\bottomrule
\end{tabularx}
\caption{Cada defecto hace fallar al menos una prueba. M11 solo lo ve la prueba unitaria, como anticipa
la gu\'ia; M5, M12 y M13 son de cableado y no tienen prueba unitaria.}
\end{table}

\subsection{Qu\'e sali\'o mal y c\'omo se encontr\'o}
\textbf{1. La unidad de forwarding se quedaba ``congelada''.} La primera versi\'on calculaba cada
salida con una funci\'on que recib\'ia solo \cod{rs} como argumento y le\'ia \cod{rdM} y
\cod{RegWriteM} directamente del m\'odulo:
\begin{lstlisting}[language=SV, numbers=none]
function automatic logic [1:0] decidir(input logic [4:0] rsE);
    if ((rsE != 5'd0) && RegWriteM && (rsE == rdM)) decidir = 2'b10; ...
endfunction
assign ForwardAE = decidir(rs1E);
\end{lstlisting}
En un \cod{assign}, el simulador solo vuelve a evaluar la funci\'on cuando cambian sus
\emph{argumentos}. En el bucle de \cod{acumulador.s}, \cod{addi x3, x3, -1} va seguido de
\cod{bne x3, x0, bucle}: \cod{rs1E} vale 3 en dos ciclos seguidos, as\'i que \cod{ForwardAE} se qued\'o
en \cod{00}, el \cod{bne} compar\'o el \cod{x3} viejo y el bucle dio una vuelta de m\'as
(\cod{x2 = 20} en vez de 15). Se encontr\'o con una traza ciclo a ciclo: en c13 se ve\'ia
\cod{rs1E = rdM = 3} y \cod{RegWriteM = 1}, pero \cod{ForwardAE = 00}. Se reescribi\'o con
\cod{always\_comb}. La prueba unitaria no lo hab\'ia detectado porque siempre cambiaba \cod{rs} al
final; ahora fija \cod{rs} primero y cambia el resto despu\'es, y con eso la versi\'on antigua da
13 fallos.

\textbf{2. Dos programas eran ciegos a ``parar sin vaciar ID/EX''.} Con ese defecto (M7), el
\cod{add} dependiente se ejecuta dos veces: la primera con la \emph{direcci\'on} del \cod{lw}
adelantada desde EX/MEM. En la primera versi\'on de \cod{acumulador.s} esa direcci\'on era 0 (sumar
0 no se nota), y en \cod{lw\_use.s} el dato~4 viv\'ia en la direcci\'on~4. El mutante pasaba todo. Se
movieron los datos a las direcciones 16, 20 y 24, y ahora M7 falla en \cod{acumulador.s}.

% =====================================================================
\section{Ciclos: segmentado frente a monociclo}
% =====================================================================

El testbench mide los ciclos de los dos procesadores con el mismo programa. En el monociclo, los
ciclos son las instrucciones ejecutadas ($N$, con CPI $=1$). En el segmentado, si la instrucci\'on
\cod{fin} entra a IF en el ciclo $c$, la anterior termina WB al final de $c+3$. Adem\'as se cuentan las
paradas ($S$) y los saltos tomados ($F$), y se verifica la f\'ormula
\[
\boxed{\;T_{\text{seg}} = N + 4 + S + 2F\;}
\]
donde 4 son los ciclos de llenado del pipeline, cada parada cuesta 1 ciclo y cada salto tomado 2.

\begin{table}[!htbp]
\centering\small
\begin{tabular}{@{}l r r r r r c@{}}
\toprule
\textbf{Programa} & \textbf{$N$ (mono)} & \textbf{$T_{\text{seg}}$} & \textbf{$S$} & \textbf{$F$} &
\textbf{CPI seg} & \textbf{$N+4+S+2F$} \\
\midrule
\cod{humo} & 4 & 8 & 0 & 0 & \num{2.00} & $4+4+0+0 = 8$ \checkmark \\
\cod{completo} & 23 & 33 & 0 & 3 & \num{1.43} & $23+4+0+6 = 33$ \checkmark \\
\cod{bytes} & 15 & 19 & 0 & 0 & \num{1.27} & $15+4+0+0 = 19$ \checkmark \\
\cod{fwd\_dist} & 10 & 14 & 0 & 0 & \num{1.40} & $10+4+0+0 = 14$ \checkmark \\
\cod{doble} & 12 & 16 & 0 & 0 & \num{1.33} & $12+4+0+0 = 16$ \checkmark \\
\cod{x0} & 8 & 12 & 0 & 0 & \num{1.50} & $8+4+0+0 = 12$ \checkmark \\
\cod{sw\_rd} & 15 & 19 & 0 & 0 & \num{1.27} & $15+4+0+0 = 19$ \checkmark \\
\cod{lui\_jal} & 13 & 21 & 0 & 2 & \num{1.62} & $13+4+0+4 = 21$ \checkmark \\
\cod{lw\_use} & 15 & 22 & 3 & 0 & \num{1.47} & $15+4+3+0 = 22$ \checkmark \\
\cod{salto} & 27 & 45 & 0 & 7 & \num{1.67} & $27+4+0+14 = 45$ \checkmark \\
\cod{acumulador} & 23 & 34 & 3 & 2 & \num{1.48} & $23+4+3+4 = 34$ \checkmark \\
\bottomrule
\end{tabular}
\caption{Ciclos medidos en simulaci\'on. La \'ultima columna se calcul\'o aparte y coincide en los once
programas.}
\end{table}

\textbf{Verificaci\'on a mano con \cod{completo.s}.} Instrucciones ejecutadas: 13 hasta el segundo
\cod{beq} (incluido); luego \cod{addi x11}, \cod{addi x12}, \cod{lui}, \cod{addi x14} y \cod{jal}
(18); \cod{add} y \cod{jalr} en \cod{duplica} (20); y \cod{addi}, \cod{sltu}, \cod{addi} al volver
(23). Saltos tomados: el segundo \cod{beq}, \cod{jal} y \cod{jalr} ($F = 3$). No hay \cod{lw}
seguido de uso ($S = 0$). Entonces $T_{\text{seg}} = 23 + 4 + 0 + 2\cdot 3 = 33$ ciclos, que coincide
con la simulaci\'on, y $\mathrm{CPI} = 33/23 = \num{1.43}$.

\textbf{En tiempo, no en ciclos.} El segmentado tarda m\'as ciclos, pero su ciclo de reloj es m\'as
corto, porque dura la etapa m\'as lenta y no el camino completo. La gu\'ia no da retardos, as\'i que solo
se puede estimar con un supuesto ideal: cinco etapas balanceadas y sin sobrecosto de los registros, es
decir, $T_{c,\text{mono}} = 5\,T_{c,\text{seg}}$. Con eso:
\[
\text{Aceleraci\'on} = \frac{N \cdot T_{c,\text{mono}}}{T_{\text{seg}}\cdot T_{c,\text{seg}}}
= \frac{23 \cdot 5}{33} \approx \num{3.48}.
\]
Para programas largos, el llenado (4 ciclos) se diluye y el CPI tiende a $1 + (S + 2F)/N$.

\begin{respuesta}[Resultado]
\cod{completo.s}: $\boxed{\text{monociclo } 23 \text{ ciclos},\ \text{segmentado } 33 \text{ ciclos},\
\mathrm{CPI} = \num{1.43}}$. Con etapas ideales, $\approx \num{3.5}\times$ m\'as r\'apido en tiempo.
\end{respuesta}

% =====================================================================
\section{Lo que cuesta ahora cada riesgo}
% =====================================================================
\begin{center}\small
\begin{tabular}{@{}l l l l@{}}
\toprule
\textbf{Riesgo} & \textbf{Antes (programador)} & \textbf{Ahora (hardware)} & \textbf{Comprobado en} \\
\midrule
Datos a distancia 1 & dos \cod{nop} & 0 ciclos, desde EX/MEM & \cod{fwd\_dist.s} \\
Datos a distancia 2 & un \cod{nop} & 0 ciclos, desde MEM/WB & \cod{fwd\_dist.s} \\
Datos a distancia $\geq 3$ & ninguno & 0 ciclos, el banco & \cod{completo.s} (M12) \\
\cod{lw} seguido de uso & dos \cod{nop} & 1 ciclo, evitable reordenando & \cod{lw\_use.s} \\
Salto tomado & dos ranuras a mano & 2 ciclos, autom\'aticos & \cod{salto.s} \\
\bottomrule
\end{tabular}
\end{center}
El binario deja de estar atado a la profundidad del pipeline: los mismos \cod{.hex} del Taller~2 corren
sin cambios y dejan el mismo estado.

% =====================================================================
\section{Consideraciones de dise\~no}
% =====================================================================
\begin{itemize}
\item \textbf{Escritura del banco en el flanco de bajada} (secci\'on 2): se conect\'o el reloj invertido a
  \cod{register\_file.sv} para que escriba en el primer medio ciclo sin modificar el m\'odulo.
\item \textbf{\cod{funct3}, \cod{ALUSrcA} y \cod{Jalr}} viajan por los registros de segmentaci\'on
  aunque la tabla de la secci\'on~2 no los mencione; sin ellos no funcionar\'ian los saltos condicionales,
  los accesos de byte y media palabra, \cod{auipc} ni \cod{jalr}.
\item \textbf{\cod{PCSrcE}} se calcula en EX con la misma l\'ogica del bloque~3 de
  \cod{control\_unit}, que se instancia en ID sin cambios (con las banderas a cero).
\item \textbf{Condici\'on de parada.} Se implement\'o tal como la define la gu\'ia, sin excluir
  \cod{rdE = x0} ni distinguir las instrucciones que no usan \cod{rs2}. En una instrucci\'on tipo~I, los
  bits 24:20 pertenecen al inmediato; si coinciden con el \cod{rd} de un \cod{lw} anterior, el pipeline
  se detiene un ciclo de m\'as. El resultado sigue siendo correcto y ninguno de los programas de prueba
  presenta ese caso. Se mantuvo as\'i porque la tabla de la secci\'on~12 es la especificaci\'on.
\item \textbf{\cod{EX/MEM} lleva \cod{PCTarget}} aunque ninguna etapa posterior lo use, para respetar la
  tabla de la secci\'on~2.
\item \textbf{Burbuja = instrucci\'on \cod{0x00000000}.} El opcode \cod{0000000} no corresponde a
  ninguna de las nueve filas de la tabla de control, as\'i que la unidad de control deja todas sus
  se\~nales en el valor por defecto (sin escritura); el borrado de IF/ID equivale a un \cod{nop}.
\item \textbf{Entorno de simulaci\'on.} Icarus Verilog 12. El c\'odigo es SystemVerilog sintetizable
  est\'andar; solo los testbenches usan referencias jer\'arquicas, igual que los del Taller~2.
\end{itemize}

\appendix
% =====================================================================
\section{Los programas de prueba}
% =====================================================================
\begin{lstlisting}[language=RVASM,title={\small\ttfamily prog/fwd\_dist.s}]
# =============================================================================
# fwd_dist.s - Dependencias a distancia 1 y 2, por el operando A y por el B
#
# Se usa sub (no add) para que el orden de los operandos importe: si A y B se
# cruzan, o si solo existe ForwardAE, el resultado cambia.
#
# Estado esperado:
#   x1 = 7   x2 = 3   x3 = 4   x4 = -1 (0xFFFFFFFF)   x5 = 4   x6 = 1
#   x7 = -1  x8 = -1
# =============================================================================
    addi  x1, x0, 7         # x1 = 7
    addi  x2, x0, 3         # x2 = 3
    sub   x3, x1, x2        # A a distancia 2 (x1), B a distancia 1 (x2): 4
    sub   x4, x2, x3        # A a distancia 2 (x2), B a distancia 1 (x3): -1
    sub   x5, x3, x0        # A a distancia 2 (x3): 4
    sub   x6, x0, x4        # B a distancia 2 (x4): 1
    sub   x7, x0, x6        # B a distancia 1 (x6): -1
    add   x8, x7, x0        # A a distancia 1 (x7): -1
    nop
    nop
fin:
    beq   x0, x0, fin
\end{lstlisting}
\begin{lstlisting}[language=RVASM,title={\small\ttfamily prog/doble.s}]
# =============================================================================
# doble.s - Doble hazard: dos escrituras seguidas al mismo registro y una
# lectura. Debe ganar EX/MEM (la mas reciente). Se prueba por A y por B.
#
# Estado esperado:
#   x5 = 70  x6 = 71  x7 = 70  x8 = 71
# Con la prioridad invertida saldria x6 = 4 y x8 = 4: plausible y equivocado.
# =============================================================================
    addi  x1, x0, 1
    addi  x2, x0, 2
    addi  x3, x0, 30
    addi  x4, x0, 40
    add   x5, x1, x2        # x5 = 3
    add   x5, x3, x4        # x5 = 70   (vuelve a escribir x5)
    add   x6, x5, x1        # lee x5 por A: 71
    add   x7, x1, x2        # x7 = 3
    add   x7, x3, x4        # x7 = 70
    add   x8, x1, x7        # lee x7 por B: 71
    nop
    nop
fin:
    beq   x0, x0, fin
\end{lstlisting}
\begin{lstlisting}[language=RVASM,title={\small\ttfamily prog/x0.s}]
# =============================================================================
# x0.s - Instrucciones que escriben en x0 seguidas de lecturas de x0
#
# Sin la condicion rs != 0, el resultado tirado a x0 se adelanta y x0 deja de
# valer cero.
#
# Estado esperado:
#   x0 = 0   x1 = 5   x2 = 5   x3 = 1   x4 = 5
# =============================================================================
    addi  x1, x0, 5         # x1 = 5
    addi  x0, x0, 77        # escribe en x0: se tira
    add   x2, x0, x1        # x0 por A a distancia 1: x2 = 5 (mal: 82)
    addi  x0, x1, 9         # otra escritura a x0 (14)
    addi  x3, x0, 1         # x0 por A a distancia 1: x3 = 1 (mal: 15)
    add   x4, x1, x0        # x0 por B a distancia 2: x4 = 5 (mal: 19)
    nop
    nop
fin:
    beq   x0, x0, fin
\end{lstlisting}
\begin{lstlisting}[language=RVASM,title={\small\ttfamily prog/sw\_rd.s}]
# =============================================================================
# sw_rd.s - Un sw (y un branch) cuyos bits 11:7 coinciden con un registro que
# se lee justo despues.
#
# "sw x10, 8(x0)" tiene instr[11:7] = 8: sin la condicion RegWrite se anuncia
# como escritura a x8 y adelanta la direccion (8) en vez del contenido de x8.
#
# Estado esperado:
#   x8 = 100   x10 = 4   x12 = 55   x11 = 100   x13 = 55   x14 = 100
#   Mem[2] (dir 8) = 4   Mem[3] (dir 12) = 4
# =============================================================================
    addi  x8,  x0, 100      # x8 = 100
    addi  x12, x0, 55       # x12 = 55
    addi  x10, x0, 4        # x10 = 4
    nop
    nop
    nop
    sw    x10, 8(x0)        # bits 11:7 = 8  -> "rd" falso x8
    add   x11, x8, x0       # x8 por A a distancia 1: 100 (mal: 8)
    sw    x10, 12(x0)       # bits 11:7 = 12 -> "rd" falso x12
    nop
    add   x13, x0, x12      # x12 por B a distancia 2: 55 (mal: 12)
    bne   x0, x0, 8         # no tomado; bits 11:7 = 8 -> "rd" falso x8
    add   x14, x0, x8       # x8 por B a distancia 1: 100 (mal: 0, el resultado del sub)
    nop
    nop
fin:
    beq   x0, x0, fin
\end{lstlisting}
\begin{lstlisting}[language=RVASM,title={\small\ttfamily prog/lui\_jal.s}]
# =============================================================================
# lui_jal.s - Instrucciones que escriben en rd algo que NO paso por la ALU
# (lui, jal, jalr) o que la ALU calcula con el PC (auipc), seguidas de un uso
# a distancia 1 y 2.
#
# Si EX/MEM solo lleva el resultado de la ALU, se adelanta basura.
#
# Estado esperado (las direcciones son las de este programa):
#   x1 = 0x12345000   x2 = 0x12345678   x3 = 0xABCDE000   x4 = 0xABCDE000
#   x5 = 0x14 (PC+4 del jal en 0x10)    x6 = 0x14
#   x7 = 0x18 (auipc en 0x18)           x8 = 0x1C
#   x9 = 0x28 (PC+4 del jalr en 0x24)   x10 = 0x28
# =============================================================================
    lui   x1, 0x12345       # 0x00  x1 = 0x12345000
    addi  x2, x1, 0x678     # 0x04  distancia 1 desde lui: 0x12345678
    lui   x3, 0xABCDE       # 0x08  x3 = 0xABCDE000
    add   x4, x0, x3        # 0x0C  distancia 1 por B: 0xABCDE000
    jal   x5, sig           # 0x10  x5 = 0x14
sig:
    addi  x6, x5, 0         # 0x14  distancia 1 desde jal: 0x14
    auipc x7, 0             # 0x18  x7 = 0x18
    addi  x8, x7, 4         # 0x1C  distancia 1 desde auipc: 0x1C
    addi  x11, x0, 0x2C     # 0x20  direccion de "ret_aqui"
    jalr  x9, 0(x11)        # 0x24  x9 = 0x28, salta a 0x2C
    addi  x20, x0, 99       # 0x28  ranura: no debe ejecutarse
ret_aqui:
    addi  x10, x9, 0        # 0x2C  distancia 1 desde jalr: 0x28
    nop
    nop
fin:
    beq   x0, x0, fin
\end{lstlisting}
\begin{lstlisting}[language=RVASM,title={\small\ttfamily prog/lw\_use.s}]
# =============================================================================
# lw_use.s - load-use por rs1 y por rs2, un load con una instruccion
# independiente antes del uso (distancia 2: no debe parar) y un load -> store.
#
# Los datos viven en direcciones distintas de su valor (21 en la 16, 4 en la
# 20): asi un uso que reciba la direccion en vez del dato se nota.
#
# Paradas esperadas: exactamente 3 (rs1, rs2 y el lw -> sw).
#
# Estado esperado:
#   x1 = 21  x2 = 4  x3 = 21  x4 = 21  x5 = 4  x6 = 4  x7 = 21  x8 = 42
#   x9 = 1   x12 = 21
#   Mem[4] = 21   Mem[5] = 4   Mem[6] = 21
# =============================================================================
    addi  x1, x0, 21
    sw    x1, 16(x0)        # store a distancia 1 del addi: el dato se adelanta
    addi  x2, x0, 4
    sw    x2, 20(x0)
    lw    x3, 16(x0)
    add   x4, x3, x0        # load-use por rs1: x4 = 21   (1 parada)
    lw    x5, 20(x0)
    add   x6, x0, x5        # load-use por rs2: x6 = 4    (1 parada)
    lw    x7, 16(x0)
    addi  x9, x0, 1         # independiente
    add   x8, x7, x7        # distancia 2: sin parada, x8 = 42
    lw    x12, 16(x0)
    sw    x12, 24(x0)       # load -> store: el dato del sw depende del lw (1 parada)
    nop
    nop
fin:
    beq   x0, x0, fin
\end{lstlisting}
\begin{lstlisting}[language=RVASM,title={\small\ttfamily prog/salto.s}]
# =============================================================================
# salto.s - Saltos tomados con instrucciones detras que escriben algo visible
# (registros y memoria). Si falta FlushD o FlushE, dejan rastro.
#
# Cubre beq, bne, blt, bge, bltu, bgeu (tomados y no tomados) y jal.
#
# Estado esperado:
#   x1 = 3  x2 = 3  x3 = 3  x4 = 7  x5 = -1  x6 = 1
#   x21 = 6 (uno por cada condicional no tomado)
#   x28 = x29 = x30 = x31 = 0       Mem[0] = Mem[1] = 0
# =============================================================================
    addi  x1, x0, 3
    addi  x2, x0, 3
    add   x3, x1, x0
    beq   x3, x2, destino   # tomado (con forwarding de x3 a distancia 1)
    addi  x30, x0, 99       # ranura 1: no debe dejar rastro
    sw    x1, 0(x0)         # ranura 2: no debe escribir memoria
destino:
    addi  x4, x0, 7
    jal   x0, sig           # salto incondicional: tambien vacia dos ranuras
    addi  x31, x0, 99
    sw    x1, 4(x0)
sig:
    addi  x5, x0, -1
    addi  x6, x0, 1
    # --- Los seis condicionales, tomados: cada ranura escribiria x28/x29 ---
    bne   x5, x6, t1
    addi  x28, x0, 1
    addi  x29, x0, 1
t1: blt   x5, x6, t2        # -1 < 1 con signo
    addi  x28, x0, 2
    addi  x29, x0, 2
t2: bge   x6, x5, t3        # 1 >= -1 con signo
    addi  x28, x0, 3
    addi  x29, x0, 3
t3: bltu  x6, x5, t4        # 1 < 0xFFFFFFFF sin signo
    addi  x28, x0, 4
    addi  x29, x0, 4
t4: bgeu  x5, x6, t5        # 0xFFFFFFFF >= 1 sin signo
    addi  x28, x0, 5
    addi  x29, x0, 5
t5:
    # --- Los seis, NO tomados: cada uno suma 1 a x21 ---
    beq   x5, x6, mal
    addi  x21, x21, 1
    bne   x1, x2, mal
    addi  x21, x21, 1
    blt   x6, x5, mal
    addi  x21, x21, 1
    bge   x5, x6, mal
    addi  x21, x21, 1
    bltu  x5, x6, mal
    addi  x21, x21, 1
    bgeu  x6, x5, mal
    addi  x21, x21, 1
    nop
    nop
fin:
    beq   x0, x0, fin
mal:
    addi  x31, x0, 66
    beq   x0, x0, fin
\end{lstlisting}
\begin{lstlisting}[language=RVASM,title={\small\ttfamily prog/acumulador.s}]
# =============================================================================
# acumulador.s - lw seguido de uso dentro de un bucle con acumulador.
#
# Si se para sin meter burbuja, el add se ejecuta dos veces por vuelta: la
# primera con la DIRECCION del lw adelantada desde EX/MEM y la segunda con el
# dato. Por eso el dato vive en la direccion 16 y no en la 0: con direccion 0
# la ejecucion repetida sumaria cero y el defecto quedaria invisible (es lo
# que paso en la primera version de este programa).
#
# Estado esperado:
#   x1 = 5   x2 = 15 (3 vueltas de +5)   x3 = 0   x4 = 5
#   x5 = 3 (una vez por vuelta)          Mem[4] = 5   Mem[5] = 15
# =============================================================================
    addi  x1, x0, 5
    sw    x1, 16(x0)        # Mem[4] (dir 16) = 5
    addi  x2, x0, 0         # acumulador
    addi  x3, x0, 3         # contador
    addi  x5, x0, 0
bucle:
    lw    x4, 16(x0)
    add   x2, x2, x4        # load-use con acumulador
    addi  x5, x5, 1         # otro acumulador
    addi  x3, x3, -1
    bne   x3, x0, bucle     # tomado 2 veces
    sw    x2, 20(x0)        # Mem[5] (dir 20) = 15
    nop
    nop
fin:
    beq   x0, x0, fin
\end{lstlisting}




\end{document}
